\documentclass[10pt,letterpaper]{article}
\usepackage[utf8]{inputenc}
\usepackage[T1]{fontenc}
\usepackage[spanish,es-noquoting,es-nodecimaldot]{babel}
\usepackage[margin=1.2cm, top=1.2cm, bottom=1.2cm, left=1.2cm, right=1.2cm]{geometry}
\usepackage{multicol}
\usepackage{listings}
\usepackage{xcolor}
\usepackage{enumitem}
\usepackage{titlesec}
\usepackage{amsmath,amssymb}

% Formato compacto de secciones para maximizar aprovechamiento de espacio
\titleformat{\section}{\large\bfseries}{\thesection.}{0.5em}{}[\titlerule]
\titleformat{\subsection}{\normalsize\bfseries}{\thesubsection.}{0.4em}{}
\titlespacing*{\section}{0pt}{8pt}{4pt}
\titlespacing*{\subsection}{0pt}{6pt}{2pt}
\setlist{nosep, leftmargin=1.3em, topsep=2pt, itemsep=2pt}

% Configuración de listings para código de gramáticas
\definecolor{codebg}{rgb}{0.97,0.97,0.97}
\lstdefinestyle{grammarstyle}{
  backgroundcolor=\color{codebg},
  basicstyle=\ttfamily\scriptsize,
  breaklines=true,
  breakatwhitespace=true,
  columns=flexible,
  keepspaces=true,
  showstringspaces=false,
  frame=single,
  rulecolor=\color{black!15},
  framerule=0.5pt,
  framesep=2pt,
  aboveskip=2pt,
  belowskip=2pt
}

\begin{document}
\pagestyle{plain}

% Encabezado técnico compacto (sin gran título ni márgenes desperdiciados)
\noindent
\begin{minipage}[t]{0.70\textwidth}
  {\Large \textbf{Gramática para Analizador Sintáctico -- Lua 5.5}} \\
  \footnotesize Lenguajes de Programación \textbullet{} Práctica 2: Analizador Sintáctico Predictivo LL(1)
\end{minipage}%
\hfill
\begin{minipage}[t]{0.28\textwidth}
  \raggedleft
  \normalsize \textbf{Jonathan Felipe López Núñez} \\
  \footnotesize Especificación
\end{minipage}

\vspace{3pt}
\hrule height 0.8pt
\vspace{6pt}

% ==============================================================================
% SECCIÓN 1: GRAMÁTICA DE GRAMATICA.TXT (SOLO LA GRAMÁTICA, SIN EXPLICACIÓN)
% ==============================================================================
\section{Gramática usada para el analizador sintáctico}

\begin{multicols}{2}
\begin{lstlisting}[style=grammarstyle]
CHUNK ::= BLOCK

BLOCK ::= STAT_LIST

STAT_LIST ::= STAT STAT_LIST
           | epsilon

STAT ::= ID_NAME PRIMARY_STAT_TAIL
       | local VAR_OR_FUNC_DECL
       | global VAR_OR_FUNC_DECL
       | function FUNC_NAME FUNC_BODY
       | while EXP do BLOCK end
       | repeat BLOCK until EXP
       | if EXP then BLOCK ELSEIF_LIST ELSE_OPT end
       | for ID_NAME FOR_TAIL
       | goto ID_NAME
       | tkn_goto ID_NAME tkn_goto
       | break
       # Segun el test Campuswire 6, se debe tener return aqui
       | return RET_TAIL
       | tkn_semicolon

ID_NAME ::= id
          | print
          | warn
          | pcall
          | xpcall
          | error

ELSEIF_LIST ::= elseif EXP then BLOCK ELSEIF_LIST
             | epsilon

ELSE_OPT ::= else BLOCK
          | epsilon

RET_TAIL ::= epsilon
          | VAL_EXP_LIST

VAR_OR_FUNC_DECL ::= function ID_NAME FUNC_BODY
                  | NAME_DECL_LIST OPT_ASSIGN

NAME_DECL_LIST ::= NAME_DECL NAME_DECL_LIST_TAIL

NAME_DECL_LIST_TAIL ::= tkn_comma NAME_DECL NAME_DECL_LIST_TAIL
                     | epsilon

NAME_DECL ::= ID_NAME NAME_DECL_TAIL

NAME_DECL_TAIL ::= tkn_opening_bra EXP tkn_closing_bra
                | epsilon

OPT_ASSIGN ::= tkn_assign VAL_EXP_LIST
            | epsilon

PRIMARY_STAT_TAIL ::= tkn_assign VAL_EXP_LIST
                   | tkn_comma VAR_LIST tkn_assign VAL_EXP_LIST
                   | tkn_colon ID_NAME CALL_TAIL
                   | tkn_opening_par ARGS_OPT tkn_closing_par CALL_STAT_TAIL
                   | tkn_period ID_NAME PRIMARY_STAT_TAIL
                   | tkn_opening_bra EXP tkn_closing_bra PRIMARY_STAT_TAIL

CALL_TAIL ::= tkn_opening_par ARGS_OPT tkn_closing_par CALL_STAT_TAIL

CALL_STAT_TAIL ::= tkn_opening_par ARGS_OPT tkn_closing_par CALL_STAT_TAIL
                | tkn_colon ID_NAME CALL_TAIL
                | tkn_period ID_NAME PRIMARY_STAT_TAIL
                | tkn_opening_bra EXP tkn_closing_bra PRIMARY_STAT_TAIL
                | epsilon

VAR_LIST ::= VAR_ITEM VAR_LIST_TAIL

VAR_LIST_TAIL ::= tkn_comma VAR_ITEM VAR_LIST_TAIL
               | epsilon

VAR_ITEM ::= ID_NAME VAR_SUFFIX

VAR_SUFFIX ::= tkn_period ID_NAME VAR_SUFFIX
            | tkn_opening_bra EXP tkn_closing_bra VAR_SUFFIX
            | epsilon

FOR_TAIL ::= tkn_assign EXP tkn_comma EXP FOR_STEP_OPT do BLOCK end
          | tkn_comma ID_NAME FOR_NAME_LIST in EXP_LIST do BLOCK end
          | in EXP_LIST do BLOCK end

FOR_STEP_OPT ::= tkn_comma EXP
              | epsilon

FOR_NAME_LIST ::= tkn_comma ID_NAME FOR_NAME_LIST
               | epsilon

FUNC_NAME ::= ID_NAME FUNC_NAME_TAIL

FUNC_NAME_TAIL ::= tkn_period ID_NAME FUNC_NAME_TAIL
                | tkn_colon ID_NAME
                | epsilon

FUNC_BODY ::= tkn_opening_par PAR_LIST_OPT tkn_closing_par BLOCK end

PAR_LIST_OPT ::= PAR_LIST
              | epsilon

PAR_LIST ::= id PAR_LIST_TAIL
          | tkn_varargs

PAR_LIST_TAIL ::= tkn_comma PAR_ITEM
               | epsilon

PAR_ITEM ::= id PAR_LIST_TAIL
          | tkn_varargs

VAL_EXP ::= function FUNC_BODY
          | EXP

VAL_EXP_LIST ::= VAL_EXP VAL_EXP_LIST_TAIL

VAL_EXP_LIST_TAIL ::= tkn_comma VAL_EXP VAL_EXP_LIST_TAIL
                   | epsilon

EXP_LIST ::= EXP EXP_LIST_TAIL

EXP_LIST_TAIL ::= tkn_comma EXP EXP_LIST_TAIL
               | epsilon

ARGS_OPT ::= VAL_EXP_LIST
          | epsilon

EXP ::= AND_EXP EXP_TAIL

EXP_TAIL ::= or AND_EXP EXP_TAIL
          | epsilon

AND_EXP ::= REL_EXP AND_EXP_TAIL

AND_EXP_TAIL ::= and REL_EXP AND_EXP_TAIL
              | epsilon

REL_EXP ::= BIT_OR_EXP REL_EXP_TAIL

REL_EXP_TAIL ::= REL_OP BIT_OR_EXP
              | epsilon

REL_OP ::= tkn_less
        | tkn_greater
        | tkn_leq
        | tkn_geq
        | tkn_neq
        | tkn_equal

BIT_OR_EXP ::= BIT_XOR_EXP BIT_OR_EXP_TAIL

BIT_OR_EXP_TAIL ::= tkn_bit_or BIT_XOR_EXP BIT_OR_EXP_TAIL
                 | epsilon

BIT_XOR_EXP ::= BIT_AND_EXP BIT_XOR_EXP_TAIL

BIT_XOR_EXP_TAIL ::= tkn_bitex_or BIT_AND_EXP BIT_XOR_EXP_TAIL
                  | epsilon

BIT_AND_EXP ::= SHIFT_EXP BIT_AND_EXP_TAIL

BIT_AND_EXP_TAIL ::= tkn_bit_and SHIFT_EXP BIT_AND_EXP_TAIL
                  | epsilon

SHIFT_EXP ::= CONCAT_EXP SHIFT_EXP_TAIL

SHIFT_EXP_TAIL ::= tkn_left_shift CONCAT_EXP SHIFT_EXP_TAIL
                | tkn_right_shift CONCAT_EXP SHIFT_EXP_TAIL
                | epsilon

CONCAT_EXP ::= ADD_EXP CONCAT_EXP_TAIL

CONCAT_EXP_TAIL ::= tkn_concat CONCAT_EXP
                 | epsilon

ADD_EXP ::= MUL_EXP ADD_EXP_TAIL

ADD_EXP_TAIL ::= tkn_plus MUL_EXP ADD_EXP_TAIL
              | tkn_minus MUL_EXP ADD_EXP_TAIL
              | epsilon

MUL_EXP ::= UNARY_EXP MUL_EXP_TAIL

MUL_EXP_TAIL ::= tkn_times UNARY_EXP MUL_EXP_TAIL
              | tkn_div UNARY_EXP MUL_EXP_TAIL
              | tkn_floor_div UNARY_EXP MUL_EXP_TAIL
              | tkn_mod UNARY_EXP MUL_EXP_TAIL
              | epsilon

UNARY_EXP ::= not UNARY_EXP
           | tkn_length UNARY_EXP
           | tkn_minus UNARY_EXP
           | POW_EXP

POW_EXP ::= ATOM POW_EXP_TAIL

POW_EXP_TAIL ::= tkn_power UNARY_EXP
              | epsilon

ATOM ::= ID_NAME SUFFIX_EXP
      | valor_numerico
      | cadena_de_caracteres
      | nil
      | true
      | false
      | TABLE_CONSTRUCTOR
      | tkn_opening_par EXP tkn_closing_par

SUFFIX_EXP ::= tkn_period ID_NAME SUFFIX_EXP
            | tkn_opening_bra EXP tkn_closing_bra SUFFIX_EXP
            | tkn_opening_par ARGS_OPT tkn_closing_par SUFFIX_EXP
            | tkn_colon ID_NAME tkn_opening_par ARGS_OPT tkn_closing_par SUFFIX_EXP
            | epsilon

TABLE_CONSTRUCTOR ::= tkn_opening_key FIELD_LIST_OPT tkn_closing_key

FIELD_LIST_OPT ::= FIELD FIELD_LIST_TAIL
                | epsilon

FIELD_LIST_TAIL ::= FIELD_SEP FIELD_OPT
                 | epsilon

FIELD_OPT ::= FIELD FIELD_LIST_TAIL
           | epsilon

FIELD_SEP ::= tkn_comma
           | tkn_semicolon

FIELD ::= tkn_opening_bra EXP tkn_closing_bra tkn_assign VAL_EXP
       | ID_NAME FIELD_ID_TAIL
       | function FUNC_BODY
       | EXP_NON_ID

FIELD_ID_TAIL ::= tkn_assign VAL_EXP
               | SUFFIX_EXP POW_EXP_TAIL MUL_EXP_TAIL ADD_EXP_TAIL SHIFT_EXP_TAIL BIT_AND_EXP_TAIL BIT_XOR_EXP_TAIL BIT_OR_EXP_TAIL REL_EXP_TAIL AND_EXP_TAIL EXP_TAIL

EXP_NON_ID ::= AND_EXP_NON_ID EXP_TAIL

AND_EXP_NON_ID ::= REL_EXP_NON_ID AND_EXP_TAIL

REL_EXP_NON_ID ::= BIT_OR_EXP_NON_ID REL_EXP_TAIL

BIT_OR_EXP_NON_ID ::= BIT_XOR_EXP_NON_ID BIT_OR_EXP_TAIL

BIT_XOR_EXP_NON_ID ::= BIT_AND_EXP_NON_ID BIT_XOR_EXP_TAIL

BIT_AND_EXP_NON_ID ::= SHIFT_EXP_NON_ID BIT_AND_EXP_TAIL

SHIFT_EXP_NON_ID ::= CONCAT_EXP_NON_ID SHIFT_EXP_TAIL

CONCAT_EXP_NON_ID ::= ADD_EXP_NON_ID CONCAT_EXP_TAIL

ADD_EXP_NON_ID ::= MUL_EXP_NON_ID ADD_EXP_TAIL

MUL_EXP_NON_ID ::= UNARY_EXP_NON_ID MUL_EXP_TAIL

UNARY_EXP_NON_ID ::= not UNARY_EXP
                  | tkn_length UNARY_EXP
                  | tkn_minus UNARY_EXP
                  | POW_EXP_NON_ID

POW_EXP_NON_ID ::= ATOM_NON_ID POW_EXP_TAIL

ATOM_NON_ID ::= valor_numerico
             | cadena_de_caracteres
             | nil
             | true
             | false
             | TABLE_CONSTRUCTOR
             | tkn_opening_par EXP tkn_closing_par
\end{lstlisting}
\end{multicols}

\vspace{6pt}

% ==============================================================================
% SECCIÓN 2: GRAMÁTICA BNF DE LUA 5.5 REFERENCE MANUAL Y ANÁLISIS DE CAMBIOS
% ==============================================================================
\section{Gramática BNF base usada para la creación de la gramática}

\subsection{Gramática BNF Oficial Extraída de Lua 5.5 Reference Manual}

A continuación se presentan todas las producciones sintácticas extraídas del manual de referencia oficial:

\begin{multicols}{2}
\begin{lstlisting}[style=grammarstyle]
chunk ::= block

block ::= {stat}

stat ::= ';'
       | varlist '=' explist
       | functioncall
       | label
       | break
       | goto Name
       | do block end
       | while exp do block end
       | repeat block until exp
       | if exp then block {elseif exp then block} [else block] end
       | for Name '=' exp ',' exp [',' exp] do block end
       | for namelist in explist do block end
       | function funcname funcbody
       | local function Name funcbody
       | global function Name funcbody
       | local attnamelist ['=' explist]
       | global attnamelist ['=' explist]
       | global [attrib] '*'
       | return [explist] [';']

label ::= '::' Name '::'

funcname ::= Name {'.' Name} [':' Name]

varlist ::= var {',' var}

var ::= Name
      | prefixexp '[' exp ']'
      | prefixexp '.' Name

namelist ::= Name {',' Name}

explist ::= exp {',' exp}

exp ::= nil | false | true | Numeral | LiteralString | '...' | functiondef
      | prefixexp | tableconstructor | exp binop exp | unop exp

prefixexp ::= var | functioncall | '(' exp ')'

functioncall ::= prefixexp args
               | prefixexp ':' Name args

args ::= '(' [explist] ')'
       | tableconstructor
       | LiteralString

functiondef ::= function funcbody

funcbody ::= '(' [parlist] ')' block end

parlist ::= namelist [',' varargparam] | varargparam

varargparam ::= '...' [Name]

tableconstructor ::= '{' [fieldlist] '}'

fieldlist ::= field {fieldsep field} [fieldsep]

field ::= '[' exp ']' '=' exp | Name '=' exp | exp

fieldsep ::= ',' | ';'

attnamelist ::= [attrib] Name [attrib] {',' Name [attrib]}

attrib ::= '<' Name '>'
\end{lstlisting}
\end{multicols}

\subsection{Cambios Realizados a la Gramática Original y Justificación Técnica}

Para adaptar la gramática formal de Lua 5.5 a un analizador sintáctico predictivo descendente \textbf{LL(1)} y cumplir con los requerimientos de UNcode, se implementaron los siguientes cambios:

\begin{enumerate}
  \item \textbf{Eliminación de recursión por la izquierda y estratificación de expresiones (Diseño LL(1)):}
  \begin{itemize}
    \item \textit{Original:} La regla \texttt{exp ::= exp binop exp | unop exp} y la recursión mutua entre \texttt{prefixexp}, \texttt{var} y \texttt{functioncall} poseen recursión por la izquierda directa e indirecta, además de ambigüedad en precedencia y asociatividad.
    \item \textit{Cambio:} Se descompuso \texttt{exp} en 12 niveles de precedencia jerárquica: \texttt{EXP} (\texttt{or}), \texttt{AND\_EXP} (\texttt{and}), \texttt{REL\_EXP} (relacionales), \texttt{BIT\_OR\_EXP} (\texttt{|}), \texttt{BIT\_XOR\_EXP} (\texttt{\textasciitilde{}}), \texttt{BIT\_AND\_EXP} (\texttt{\&}), \texttt{SHIFT\_EXP} (\texttt{<{}<}, \texttt{>{}>}), \texttt{CONCAT\_EXP} (\texttt{..}), \texttt{ADD\_EXP} (\texttt{+}, \texttt{-}), \texttt{MUL\_EXP} (\texttt{*}, \texttt{/}, \texttt{//}, \texttt{\%}), \texttt{UNARY\_EXP} (\texttt{not}, \texttt{\#}, \texttt{-}) y \texttt{POW\_EXP} (\texttt{\textasciicircum{}}).
    \item \textit{Asociatividad:} Para operadores asociativos a derecha (\texttt{..} y \texttt{\textasciicircum{}}), la recursión gramatical se mantiene a la derecha; para asociatividad a izquierda, se introdujeron no terminales auxiliares de cola (\texttt{*\_TAIL}) con derivación a \texttt{epsilon}.
  \end{itemize}

  \item \textbf{Factorización del conflicto en sentencias: Asignación vs Llamada a Función:}
  \begin{itemize}
    \item \textit{Original:} \texttt{stat ::= varlist '=' explist | functioncall}. Ambas sentencias inician mediante \texttt{prefixexp} (típicamente un identificador \texttt{Name}), lo que genera un conflicto First/First irresoluble con un solo token de anticipación.
    \item \textit{Cambio:} Se unificaron bajo la regla factorizada \texttt{STAT ::= ID\_NAME PRIMARY\_STAT\_TAIL}. En \texttt{PRIMARY\_STAT\_TAIL} se procesan de forma uniforme accesos a campos (\texttt{.}), índices (\texttt{[...]}), llamadas (\texttt{(...)}) y métodos (\texttt{:}), bifurcando deterministamente entre llamada consumida o sentencia de asignación (\texttt{tkn\_assign} o continuación de lista con \texttt{tkn\_comma}).
  \end{itemize}

  \item \textbf{Exclusividad del token \texttt{do} para estructuras de ciclo:}
  \begin{itemize}
    \item \textit{Original:} \texttt{stat ::= do block end}.
    \item \textit{Cambio:} Se eliminó la sentencia de bloque explícito independiente. El token \texttt{do} queda reservado única y exclusivamente como delimitador en los ciclos \texttt{while ... do} y \texttt{for ... do}. Por tanto, \texttt{do} no pertenece a $\text{First}(\texttt{STAT})$ ni es esperado al inicio de una sentencia.
  \end{itemize}

  \item \textbf{Restricción de funciones anónimas en ciclo \texttt{for} numérico:}
  \begin{itemize}
    \item \textit{Original:} \texttt{stat ::= for Name '=' exp ',' exp [',' exp] do block end}.
    \item \textit{Cambio:} En la expresión inicial tras el signo \texttt{=} en la cabecera del ciclo numérico, se restringe la posibilidad de definir closures anónimos (\texttt{function funcbody}). Por ello se evalúa directamente con \texttt{EXP} (cuyo First no incluye \texttt{function}). En asignaciones estándar o retornos sí se admite función anónima mediante \texttt{VAL\_EXP ::= function FUNC\_BODY | EXP}.
  \end{itemize}

  \item \textbf{Unificación y simplificación de declaraciones \texttt{local} y \texttt{global}:}
  \begin{itemize}
    \item \textit{Original:} \texttt{stat ::= local attnamelist ['=' explist] | global attnamelist ['=' explist]} \\
    \hspace*{1.2em} $\mid$ \texttt{global [attrib] '*'}, soportando atributos como \texttt{<const>} o \texttt{<close>}.
    \item \textit{Cambio:} Se eliminaron los atributos y el comodín \texttt{global *}, unificando ambas sentencias bajo \texttt{VAR\_OR\_FUNC\_DECL}. Tras \texttt{local} o \texttt{global} únicamente se permite una definición de función (\texttt{function ID\_NAME FUNC\_BODY}) o una lista de nombres de variables con inicialización opcional (\texttt{NAME\_DECL\_LIST OPT\_ASSIGN}).
  \end{itemize}

  \item \textbf{Resolución de ambigüedad en constructores de tabla:}
  \begin{itemize}
    \item \textit{Original:} \texttt{field ::= '[' exp ']' '=' exp | Name '=' exp | exp}.
    \item \textit{Cambio:} Cuando un campo inicia con un identificador, existe un conflicto First/First entre la asignación \texttt{Name = exp} y una expresión que comience por dicho nombre (ej. \texttt{x + 1} o \texttt{f(x)}). Se resolvió factorizando mediante \texttt{FIELD ::= ID\_NAME FIELD\_ID\_TAIL}, distinguiendo de inmediato si continúa una asignación (\texttt{tkn\_assign}) o la evaluación de una expresión general (\texttt{SUFFIX\_EXP ...}). Para expresiones que inician con otros términos se utiliza \texttt{EXP\_NON\_ID}.
  \end{itemize}

  \item \textbf{Validación estricta de parámetros en funciones y varargs:}
  \begin{itemize}
    \item \textit{Original:} \texttt{parlist ::= namelist [',' varargparam] | varargparam}, con \texttt{varargparam ::= '...' [Name]}.
    \item \textit{Cambio:} Se eliminó el nombre opcional en varargs (\texttt{... [Name]} $\rightarrow$ \texttt{tkn\_varargs}). Además, en \texttt{PAR\_LIST}, toda coma debe ir obligatoriamente seguida de un nuevo parámetro o \texttt{tkn\_varargs}, no admitiéndose comas flotantes finales.
  \end{itemize}

  \item \textbf{Reutilización de funciones nativas como identificadores:}
  \begin{itemize}
    \item \textit{Original:} El manual utiliza únicamente el no terminal abstracto \texttt{Name}.
    \item \textit{Cambio:} El analizador léxico reconoce nombres de funciones del sistema (\texttt{print}, \texttt{warn}, \texttt{pcall}, \texttt{xpcall} y \texttt{error}) como tokens con nombre propio. Como en Lua es válido redefinirlos o utilizarlos como variables y nombres de funciones propias, se definió la producción:
    \[\texttt{ID\_NAME ::= id | print | warn | pcall | xpcall | error}\]
    permitiendo su uso en cualquier posición sintáctica donde se espere un identificador.
  \end{itemize}

  \item \textbf{Adaptación estricta a los tokens del Analizador Léxico y operaciones Bitwise binarias:}
  \begin{itemize}
    \item \textit{Cambio:} Se mapearon literales y operadores a los tokens exactos del analizador léxico: literales (\texttt{valor\_numerico}, \texttt{cadena\_de\_\allowbreak caracteres}, \texttt{nil}, \texttt{true}, \texttt{false}) y operadores con prefijo \texttt{tkn\_*} (\texttt{tkn\_assign}, \texttt{tkn\_semicolon}, \texttt{tkn\_plus}, etc.).
    \item \textit{Operaciones bitwise:} Conforme a las instrucciones del proyecto, las operaciones a nivel de bits se definen exclusivamente como binarias entre dos operandos (\texttt{tkn\_bit\_and}, \texttt{tkn\_bit\_or}, \texttt{tkn\_bitex\_or}, \texttt{tkn\_left\_shift}, \texttt{tkn\_right\_shift}).
  \end{itemize}

  \item \textbf{Sentencia \texttt{return} permisiva dentro de bloques:}
  \begin{itemize}
    \item \textit{Original:} En Lua estándar, \texttt{return} se define típicamente como la última sentencia del bloque (\texttt{laststat}).
    \item \textit{Cambio:} Se integró \texttt{return RET\_TAIL} como una alternativa regular dentro de \texttt{STAT}, permitiendo retornos en cualquier punto de la secuencia de sentencias (\texttt{STAT\_LIST}) de acuerdo con las pruebas del curso.
  \end{itemize}
\end{enumerate}

\end{document}
